\documentclass[12pt,a4paper]{article}
\usepackage[utf8]{inputenc}
\usepackage[utf8]{vietnam}
\usepackage{amsmath, amssymb}
\usepackage{geometry}
\geometry{a4paper, left=2cm, right=2cm, top=2cm, bottom=2cm}
\usepackage{tcolorbox}
\usepackage{xcolor}
\usepackage{enumitem}
\usepackage{tabularx}
\usepackage{mathrsfs}

\definecolor{myblue}{RGB}{0, 102, 204}
\definecolor{myred}{RGB}{204, 0, 0}
\definecolor{mygreen}{RGB}{0, 153, 76}
\definecolor{myorange}{RGB}{255, 128, 0}

\begin{document}

\begin{center}
    \Large \textbf{\textcolor{myred}{DẠNG TOÁN [T1]: KHÁI NIỆM VECTƠ TRONG KHÔNG GIAN}} \\
    \vspace{0.2cm}
    \large \textit{(Mức độ: Nhận biết - Dành cho học sinh mất gốc)}
\end{center}

\vspace{0.5cm}

% ==========================================
% KHỐI 1: DẤU HIỆU NHẬN DIỆN
% ==========================================
\begin{tcolorbox}[colback=blue!5!white, colframe=myblue, title=\textbf{KHỐI 1: DẤU HIỆU NHẬN DIỆN DẠNG TOÁN (Nhìn phát biết ngay)}]
    \textbf{1. Từ khóa trong đề bài:}
    \begin{itemize}[leftmargin=*, label=+]
        \item Các chữ: "Vectơ", "điểm đầu", "điểm cuối", "giá của vectơ", "phương", "hướng", "độ dài".
        \item Các cụm từ hỏi: "Có bao nhiêu vectơ...", "Vectơ nào bằng vectơ...", "Vectơ nào là vectơ đối...", "Mệnh đề nào đúng/sai?".
    \end{itemize}
    
    \textbf{2. Dấu hiệu hình vẽ:} 
    Hình không gian quen thuộc như \textbf{hình tứ diện, hình hộp, hình lập phương, hình chóp, hình lăng trụ}.
    
    \textbf{3. Các biến thể thường gặp:}
    \begin{itemize}[leftmargin=*, label=+]
        \item \textit{Biến thể 1:} Nhận biết vectơ qua hình vẽ thực tế (lực kéo, chuyển động...).
        \item \textit{Biến thể 2:} Đếm số vectơ tạo thành từ các đỉnh của hình không gian.
        \item \textit{Biến thể 3:} Xác định hai vectơ bằng nhau hoặc vectơ đối của một vectơ cho trước.
    \end{itemize}
\end{tcolorbox}

% ==========================================
% KHỐI 2: QUY TRÌNH 3 BƯỚC GIẢI CHUẨN
% ==========================================
\begin{tcolorbox}[colback=green!5!white, colframe=mygreen, title=\textbf{KHỐI 2: QUY TRÌNH 3 BƯỚC GIẢI CHUẨN (Step-by-Step Checklist)}]
    \textbf{BƯỚC 1: Xác định gốc và ngọn (Điểm đầu - Điểm cuối)}
    \begin{itemize}[leftmargin=*, label=-]
        \item Vẽ hình không gian ra nháp (hoặc hình dung trong đầu).
        \item Xác định rõ đâu là điểm đầu (gốc), đâu là điểm cuối (ngọn) để tạo thành vectơ \(\overrightarrow{AB}\). \textit{(Lưu ý: \(\overrightarrow{AB} \neq \overrightarrow{BA}\))}.
    \end{itemize}
    
    \textbf{BƯỚC 2: So sánh Phương - Hướng - Độ dài (Dựa vào tính chất hình học)}
    \begin{itemize}[leftmargin=*, label=-]
        \item \textbf{Độ dài:} \(|\overrightarrow{AB}| = AB\) (chính là độ dài đoạn thẳng nối hai điểm \(A\) và \(B\)).
        \item \textbf{Giá của vectơ:} Là đường thẳng đi qua điểm đầu và điểm cuối. (Hai giá song song hoặc trùng nhau $\Rightarrow$ Cùng phương).
        \item \textbf{Cùng hướng:} Cùng phương và mũi tên chỉ về cùng một phía.
        \item \textbf{Vectơ bằng nhau:} \(\vec{a} = \vec{b}\) khi và chỉ khi chúng \textbf{cùng hướng} và \textbf{cùng độ dài}.
        \item \textbf{Vectơ đối:} \(-\vec{a}\) là vectơ có \textbf{ngược hướng} và \textbf{cùng độ dài} với \(\vec{a}\).
    \end{itemize}
    
    \textbf{BƯỚC 3: Kết luận và chốt đáp án}
    \begin{itemize}[leftmargin=*, label=-]
        \item Dựa vào kết quả so sánh ở Bước 2 để trả lời đúng câu hỏi (Đếm số lượng/Chọn đáp án đúng/Giải thích).
    \end{itemize}
\end{tcolorbox}

% ==========================================
% KHỐI 3: BÍ QUYẾT CASIO & MẸO LOẠI TRỪ
% ==========================================
\begin{tcolorbox}[colback=orange!5!white, colframe=myorange, title=\textbf{KHỐI 3: BÍ QUYẾT BẤM MÁY TÍNH CASIO \& MẸO LOẠI TRỪ TRẮC NGHIỆM}]
    \textbf{1. Sử dụng Casio fx-580VNX / fx-880BTG (Nếu bài toán có số liệu tọa độ):}
    \begin{itemize}[leftmargin=*, label=+]
        \item \textbf{MENU 5 (Vectơ):} Nhập tọa độ các vectơ vào máy.
        \item \textbf{Tính độ dài vectơ:} Bấm \textbf{SHIFT} $\rightarrow$ \textbf{HYP} (Abs) $\rightarrow$ Gọi vectơ \(A\) $\rightarrow$ \textbf{=}. Kết quả chính là độ dài cần tìm.
        \item \textbf{Kiểm tra 2 vectơ bằng nhau:} Nhập hiệu \(\vec{a} - \vec{b}\). Nếu kết quả bằng \(\vec{0} (0;0;0)\) thì 2 vectơ bằng nhau.
    \end{itemize}
    
    \textbf{2. Mẹo loại trừ đáp án trắc nghiệm trong 5 giây:}
    \begin{itemize}[leftmargin=*, label=+]
        \item \textbf{Loại ngay đáp án sai về hướng:} Ví dụ tìm vectơ bằng \(\overrightarrow{AB}\), đáp án nào có điểm đầu là \(B\) (như \(\overrightarrow{BA}\) hoặc \(\overrightarrow{BC}\)) là sai luôn, vì không cùng hướng.
        \item \textbf{Loại đáp án sai về độ dài:} Trong hình hộp, các cạnh bên song song và bằng nhau, nhưng đường chéo thì dài hơn cạnh $\rightarrow$ Đừng bao giờ chọn đáp án nói vectơ cạnh bằng vectơ đường chéo.
        \item \textbf{Lưu ý "bẫy" vectơ-không:} Vectơ-không \(\vec{0}\) cùng phương, cùng hướng với \textit{mọi} vectơ nhưng độ dài bằng 0.
    \end{itemize}
\end{tcolorbox}

% ==========================================
% KHỐI 4: CÁC VÍ DỤ MINH HỌA MẪU
% ==========================================
\begin{tcolorbox}[colback=white!95!black, colframe=myred, title=\textbf{KHỐI 4: CÁC VÍ DỤ MINH HỌA MẪU (Phủ hết các biến thể)}]
    
    \textbf{VÍ DỤ 1: (Biến thể 1 - Nhận biết qua hình vẽ thực tế)} \\
    \textit{(Trích HĐ1 trang 46 - SGK Toán 12 KNTT)} \\
    Trong Hình 2.2, lực căng dây (được tạo ra bởi sức nặng của kiện hàng) được thể hiện bởi các đoạn thẳng có mũi tên màu đỏ. \\
    a) Các đoạn thẳng này cho biết gì về hướng và độ lớn của các lực căng dây? \\
    b) Các đoạn thẳng này có cùng nằm trong một mặt phẳng không?
    
    \vspace{0.2cm}
    \textbf{\textcolor{myblue}{Lời giải chi tiết:}}
    
    \textbf{[Bước 1] Xác định gốc và ngọn:} 
    Các đoạn thẳng có mũi tên (vectơ) được vẽ xuất phát từ móc cẩu hướng xuống kiện hàng (hoặc ngược lại).
    
    \textbf{[Bước 2] So sánh Phương - Hướng - Độ dài:}
    \begin{itemize}[leftmargin=*, label=-]
        \item \textit{Về hướng:} Mũi tên chỉ lên trên (về phía móc cẩu). Đây chính là hướng của các lực căng dây.
        \item \textit{Về độ lớn:} Độ dài đoạn thẳng (tính theo tỉ lệ xích) thể hiện độ lớn của các lực. Đoạn thẳng càng dài, lực càng lớn.
    \end{itemize}
    
    \textbf{[Bước 3] Kết luận:}
    \begin{itemize}[leftmargin=*, label=-]
        \item a) Đoạn thẳng cho biết hướng (lên trên) và độ lớn (độ dài đoạn thẳng) của lực căng dây.
        \item b) Các đoạn thẳng này tạo thành một hình chóp cụt/một chùm, chúng \textbf{không cùng nằm trong một mặt phẳng}.
    \end{itemize}
    
    \noindent\textbf{[💡 Mẹo \& Góc nhìn nhanh]:} Với các bài toán thực tế, vectơ biểu diễn đại lượng có hướng. Chiều mũi tên là hướng, độ dài nét vẽ là độ lớn. Không cần tính toán phức tạp, hãy nhìn vào hình vẽ không gian 3D để thấy chúng không đồng phẳng.

    \vspace{0.5cm}
    \hrule
    \vspace{0.5cm}

    \textbf{VÍ DỤ 2: (Biến thể 3 - Xác định vectơ bằng nhau, cùng phương, cùng hướng)} \\
    \textit{(Trích HĐ2 trang 47 - SGK Toán 12 KNTT)} \\
    Cho hình hộp \(ABCD.A'B'C'D'\) (H.2.7). \\
    a) So sánh độ dài hai vectơ \(\overrightarrow{AB}\) và \(\overrightarrow{D'C'}\). \\
    b) Nhận xét về giá của hai vectơ \(\overrightarrow{AB}\) và \(\overrightarrow{D'C'}\). \\
    c) Hai vectơ \(\overrightarrow{AB}\) và \(\overrightarrow{D'C'}\) có cùng phương không? Có cùng hướng không?
    
    \vspace{0.2cm}
    \textbf{\textcolor{myblue}{Lời giải chi tiết:}}
    
    \textbf{[Bước 1] Xác định gốc và ngọn:} 
    Vectơ \(\overrightarrow{AB}\) (A là gốc, B là ngọn); Vectơ \(\overrightarrow{D'C'}\) (D' là gốc, C' là ngọn).
    
    \textbf{[Bước 2] So sánh Phương - Hướng - Độ dài:}
    \begin{itemize}[leftmargin=*, label=-]
        \item \textit{Độ dài:} Hình hộp có các mặt là hình bình hành. Mặt đáy \(ABCD\) là hình bình hành nên \(AB = DC\). Mặt bên \(DCC'D'\) là hình bình hành nên \(DC = D'C'\). Từ đó suy ra \(AB = D'C'\) hay \(|\overrightarrow{AB}| = |\overrightarrow{D'C'}|\).
        \item \textit{Giá:} \(AB \parallel DC\) và \(DC \parallel D'C'\) nên \(AB \parallel D'C'\). Vậy giá của hai vectơ này song song với nhau.
        \item \textit{Hướng:} Hai vectơ cùng phương (vì giá song song). Chiều mũi tên từ \(A \to B\) và \(D' \to C'\) cùng chỉ về một phía (sang phải).
    \end{itemize}
    
    \textbf{[Bước 3] Kết luận:}
    \begin{itemize}[leftmargin=*, label=-]
        \item a) \(|\overrightarrow{AB}| = |\overrightarrow{D'C'}|\).
        \item b) Giá của chúng song song với nhau.
        \item c) Hai vectơ cùng phương và cùng hướng.
    \end{itemize}
    
    \noindent\textbf{[💡 Mẹo \& Góc nhìn nhanh]:} Trong hình hộp/hình lăng trụ, các cạnh song song và bằng nhau thì tạo ra các vectơ bằng nhau. \textbf{Chốt luôn:} \(\overrightarrow{AB} = \overrightarrow{DC} = \overrightarrow{D'C'}\) (dựa vào quy tắc hình bình hành).

    \vspace{0.5cm}
    \hrule
    \vspace{0.5cm}

    \textbf{VÍ DỤ 3: (Biến thể 3 - Dựng điểm thỏa mãn đẳng thức vectơ)} \\
    \textit{(Trích Ví dụ 2 trang 47 - SGK Toán 12 KNTT)} \\
    Cho hình lăng trụ \(ABC.A'B'C'\) (H.2.8). \\
    a) Trong ba vectơ \(\overrightarrow{BC}, \overrightarrow{CC'}\) và \(\overrightarrow{B'B}\), vectơ nào bằng vectơ \(\overrightarrow{AA'}\)? Giải thích vì sao. \\
    b) Gọi \(M\) là trung điểm của cạnh \(BC\). Xác định điểm \(M'\) sao cho \(\overrightarrow{MM'} = \overrightarrow{AA'}\).
    
    \vspace{0.2cm}
    \textbf{\textcolor{myblue}{Lời giải chi tiết:}}
    
    \textbf{[Bước 1] Xác định gốc và ngọn:} 
    Xác định các vectơ đã cho trên hình lăng trụ.
    
    \textbf{[Bước 2] So sánh Phương - Hướng - Độ dài:}
    \begin{itemize}[leftmargin=*, label=-]
        \item \textit{Đối với câu a:} Hình lăng trụ có các cạnh bên song song và bằng nhau: \(AA' \parallel CC'\) và \(AA' = CC'\). Hai vectơ \(\overrightarrow{AA'}\) và \(\overrightarrow{CC'}\) có cùng độ dài, cùng hướng (từ mặt đáy lên mặt trên) nên \(\overrightarrow{AA'} = \overrightarrow{CC'}\). Các vectơ còn lại không cùng phương hoặc ngược hướng.
        \item \textit{Đối với câu b:} Vì \(\overrightarrow{MM'} = \overrightarrow{AA'}\), nên \(MM' \parallel AA'\) và \(MM' = AA'\). Tứ giác \(BCC'B'\) là hình bình hành, \(M\) là trung điểm \(BC\) \(\Rightarrow\) Để có \(\overrightarrow{MM'}\) cùng hướng và cùng độ dài với \(\overrightarrow{BB'}\) (vì \(\overrightarrow{BB'} = \overrightarrow{AA'}\)), thì \(M'\) phải là trung điểm của \(B'C'\).
    \end{itemize}
    
    \textbf{[Bước 3] Kết luận:}
    \begin{itemize}[leftmargin=*, label=-]
        \item a) \(\overrightarrow{CC'} = \overrightarrow{AA'}\).
        \item b) \(M'\) là trung điểm của cạnh \(B'C'\).
    \end{itemize}
    
    \noindent\textbf{[💡 Mẹo \& Góc nhìn nhanh]:} Muốn tìm vectơ bằng \(\overrightarrow{AA'}\), hãy tìm một đoạn thẳng khác \textbf{song song và cùng chiều dài} với \(AA'\) trên hình. Trong lăng trụ đứng, đó chính là các cạnh bên khác (\(BB', CC'\)).
\end{tcolorbox}

% ==========================================
% KHỐI 5: BÀI TẬP RÈN LUYỆN TƯƠNG TỰ
% ==========================================
\begin{tcolorbox}[colback=white, colframe=black, title=\textbf{KHỐI 5: BÀI TẬP RÈN LUYỆN TƯƠNG TỰ (Lặp đi lặp lại để khắc sâu)}]
    
    \textbf{NHÓM 1: Tương tự Ví dụ 1 (Nhận biết vectơ cơ bản, đếm số lượng)}
    
    \textbf{Bài 1.1} \textit{(Luyện tập 1 trang 47 - SGK):} Cho hình lập phương \(ABCD.A'B'C'D'\) (H.2.6). Trong các vectơ \(\overrightarrow{AC}, \overrightarrow{AD}, \overrightarrow{AD'}\): \\
    a) Hai vectơ nào có giá cùng nằm trong mặt phẳng \((ABCD)\)? \\
    b) Hai vectơ nào có cùng độ dài?
    
    \textbf{Bài 1.2} \textit{(Trắc nghiệm Câu 1 - Phần C):} Cho tứ diện \(ABCD\). Hỏi có bao nhiêu vectơ khác vectơ \(\vec{0}\) mà mỗi vectơ có điểm đầu, điểm cuối là hai đỉnh của tứ diện \(ABCD\)?
    \begin{tcolorbox}[colback=white, colframe=white, top=0pt, bottom=0pt, left=0pt]
        \textbf{A.} 12. \hfill \textbf{B.} 4. \hfill \textbf{C.} 10. \hfill \textbf{D.} 8.
    \end{tcolorbox}
    
    \vspace{0.3cm}
    \textbf{NHÓM 2: Tương tự Ví dụ 2 (Xác định vectơ bằng nhau, mệnh đề đúng/sai)}
    
    \textbf{Bài 2.1} \textit{(Trắc nghiệm Câu 15 - Phần C):} Cho hình hộp chữ nhật \(ABCD.A'B'C'D'\). Khi đó, vectơ bằng vectơ \(\overrightarrow{AB}\) là vectơ nào dưới đây?
    \begin{tcolorbox}[colback=white, colframe=white, top=0pt, bottom=0pt, left=0pt]
        \textbf{A.} \(\overrightarrow{D'C'}\). \hfill \textbf{B.} \(\overrightarrow{BA}\). \hfill \textbf{C.} \(\overrightarrow{CD}\). \hfill \textbf{D.} \(\overrightarrow{B'A'}\).
    \end{tcolorbox}
    
    \textbf{Bài 2.2} \textit{(Trắc nghiệm đúng/sai Câu 3 - Phần D):} Cho hình lập phương \(ABCD.A'B'C'D'\). Các mệnh đề sau đúng hay sai?
    \begin{enumerate}[label=\alph*)]
        \item \(\overrightarrow{DC'} = \overrightarrow{DC} + \overrightarrow{DD'}\).
        \item \(|\overrightarrow{AB}| = |\overrightarrow{C'D'}|\).
        \item \(\overrightarrow{A'B'} = \overrightarrow{CD}\).
        \item \(\overrightarrow{AB} + \overrightarrow{AD} + \overrightarrow{AA'} = \overrightarrow{AC'}\).
    \end{enumerate}
    
    \textbf{Bài 2.3} \textit{(Trắc nghiệm Câu 2 - Phần C):} Cho hình lập phương \(ABCD.A'B'C'D'\). Mệnh đề nào sau đây \textbf{sai}?
    \begin{tcolorbox}[colback=white, colframe=white, top=0pt, bottom=0pt, left=0pt]
        \textbf{A.} \(\overrightarrow{AB} + \overrightarrow{AD} + \overrightarrow{AA'} = \overrightarrow{AC'}\). \hfill \textbf{B.} \(\overrightarrow{AC} = \overrightarrow{AB} + \overrightarrow{AD}\). \\
        \textbf{C.} \(|\overrightarrow{AB}| = |\overrightarrow{CD}|\). \hfill \textbf{D.} \(\overrightarrow{AB} = \overrightarrow{CD}\).
    \end{tcolorbox}
    
    \vspace{0.3cm}
    \textbf{NHÓM 3: Tương tự Ví dụ 3 (Xác định điểm thỏa mãn đẳng thức vectơ)}
    
    \textbf{Bài 3.1} \textit{(Luyện tập 2 trang 48 - SGK):} Cho hình chóp \(S.ABCD\) có đáy \(ABCD\) là hình bình hành. \\
    a) Trong ba vectơ \(\overrightarrow{SC}, \overrightarrow{AD}\) và \(\overrightarrow{DC}\), vectơ nào bằng vectơ \(\overrightarrow{AB}\)? \\
    b) Gọi \(M\) là một điểm thuộc cạnh \(AD\). Xác định điểm \(N\) sao cho \(\overrightarrow{MN} = \overrightarrow{AB}\).
    
    \textbf{Bài 3.2} \textit{(Vận dụng 1 trang 48 - SGK):} Một tòa nhà có chiều cao của các tầng là như nhau. Một chiếc thang máy di chuyển từ tầng 15 lên tầng 22 của tòa nhà, sau đó di chuyển từ tầng 22 lên tầng 29. Các vectơ biểu diễn độ dịch chuyển của thang máy trong hai lần di chuyển đó có bằng nhau không? Giải thích vì sao.
    
    \vspace{0.5cm}
    \hrule
    \vspace{0.3cm}
    
    \textbf{\textcolor{myred}{BẢNG ĐÁP ÁN VÀ HƯỚNG DẪN VẮN TẮT}}
    
    \begin{tabularx}{\textwidth}{|c|X|c|}
        \hline
        \textbf{Bài} & \textbf{Hướng dẫn giải vắn tắt} & \textbf{Đáp án} \\
        \hline
        1.1 & a) \(\overrightarrow{AC}, \overrightarrow{AD}\) có giá nằm trong \((ABCD)\). \newline b) \(|\overrightarrow{AC}| = |\overrightarrow{AD'}|\) (đều là đường chéo hình vuông cạnh a). & a) AC, AD \newline b) AC, AD' \\
        \hline
        1.2 & Tứ diện có 4 đỉnh. Số vectơ tạo thành từ 2 đỉnh bất kỳ (có phân biệt thứ tự gốc-ngọn) là \(4 \times 3 = 12\). & \textbf{A} \\
        \hline
        2.1 & Trong hình hộp chữ nhật, các cạnh đối song song và bằng nhau. \(\overrightarrow{AB}\) cùng hướng và cùng độ dài với \(\overrightarrow{D'C'}\). & \textbf{A} \\
        \hline
        2.2 & a) Đúng (quy tắc hình bình hành). \newline b) Đúng (cạnh hình vuông bằng nhau). \newline c) Sai (\(\overrightarrow{A'B'}\) và \(\overrightarrow{CD}\) ngược hướng). \newline d) Đúng (quy tắc hình hộp). & a) Đ \newline b) Đ \newline c) S \newline d) Đ \\
        \hline
        2.3 & D sai vì \(\overrightarrow{AB}\) và \(\overrightarrow{CD}\) ngược hướng (dù cùng độ dài). Các câu A, B, C đều đúng theo quy tắc. & \textbf{D} \\
        \hline
        3.1 & a) \(ABCD\) là hbh nên \(\overrightarrow{AB} = \overrightarrow{DC}\). \newline b) Qua \(M\) kẻ đường thẳng song song với \(AB\), cắt \(BC\) tại \(N\). Khi đó \(ABNM\) là hbh, suy ra \(\overrightarrow{MN} = \overrightarrow{AB}\). & a) \(\overrightarrow{DC}\) \newline b) \(N\) thuộc \(BC\) sao cho \(ABNM\) là hbh. \\
        \hline
        3.2 & \textbf{Có bằng nhau.} Vì thang máy đều di chuyển từ tầng thấp lên tầng cao (cùng hướng), mỗi lần đều dịch chuyển 7 tầng (cùng độ dài). & Bằng nhau \\
        \hline
    \end{tabularx}
\end{tcolorbox}

\end{document}